\section{The reverse count: sixteen patterns}\label{sec:inversa}

To determine the value of the limit, we need to evaluate the denominator in the double-counting identity.
Viewed from the layer with
one more queen, a repair is identified by the pair of queens it leaves on the board. The reverse counter
therefore tracks pairs, and
its normalization is of order $n^2$.

Throughout this section $0\le h\le K$ (including layer zero), the process is $\QQ_h$, and $w$ is a bounded
continuous function on
$[0,1]^2$ evaluated at the centers of the root squares. We study $I_w$ on the final placement.

\subsection{Patterns and an explicit compensator}

Write a compatible ordered pair as $e_1=(r,j)$, $e_2=(i,c)$. The operation inserts $e^\star=(i,j)$ and
creates the root $o=(r,c)$. The
eight lines of the pair are distinct, and the two diagonals $d_1,d_2$ of $e^\star$ are distinct from them.
We sum only over compatible
pairs. For $L\subseteq\{1,2\}$ and $J\subseteq\{1,2\}$ define \[
 Z_{L,J}=\prod_{g\in L}O_{e_g}\prod_{g\notin L}\chi_{e_g}\prod_{i\in J}\chi_{d_i}\prod_{i\notin J}V_{d_i}.
\] There are sixteen patterns per pair. Let $\ell=|L|$. For $z_t\ge\zeta$, with $\alpha=(Nz_t)^{-1}$,
$\beta_i=p_{d_i}/(Nz_t)$, and
$k_g=\pi_{e_g}/(Nz_t^4)$, the ideal generator is \[
 \cL Z_{L,J}=-\bigl[4(2-\ell)+|J|\bigr]\alpha\,Z_{L,J}-\sum_{i\notin J}\beta_i\,Z_{L,J\cup\{i\}}
 +\sum_{g\in L}k_g\,Z_{L\setminus\{g\},J}.
\] By the geometry above, a birth does not erase the other marks. The real-piece part of the death of a mark
$V$ requires the active mark and
therefore enlarges $J$. In the region~\eqref{eq:regularidad}, with tolerance $\eps n$, the remainder
relative to the true generator
satisfies \[
 |R_{L,J}|\le C_\zeta\Bigl[(\eps/n+n^{-2})Z_{L,J}+(\eps/n^2+n^{-3})\sum_{g\in L}Z_{L\setminus\{g\},J}\Bigr];
\] the bound $O(n^{-2})$ on double deletions covers both real pieces, by codegree one, and virtual pieces,
by their Cartesian structure.

Fix a deterministic horizon $T\le T_\zeta$ and set $\Delta y=(T-t)/N$, $\theta_g=\pi_{e_g}\Delta y/z_t^4$,
and $\xi_i=p_{d_i}\Delta y/z_t$.

\begin{lema}[Compensator]\label{lem:compensador} The function \[
 \Phi^{(e_1,e_2)}_t=\bigl(O_{e_1}+\theta_1\chi_{e_1}\bigr)\bigl(O_{e_2}+\theta_2\chi_{e_2}\bigr)
 \bigl(V_{d_1}-\xi_1\chi_{d_1}\bigr)\bigl(V_{d_2}-\xi_2\chi_{d_2}\bigr)
\] solves $\partial_t\Phi+\cL\Phi=0$ with terminal value $O_{e_1}O_{e_2}V_{d_1}V_{d_2}$ at $t=T$. Its
coefficients
$c_{L,J}=\prod_{g\notin L}\theta_g\prod_{i\in J}(-\xi_i)$ in the basis $Z_{L,J}$ satisfy
$|c_{L,J}|\le C_\zeta n^{-(2-\ell)}$,
$|\partial_tc_{L,J}|\le C_\zeta n^{-(3-\ell)}$, and $|\partial_t^2c_{L,J}|\le C_\zeta n^{-(4-\ell)}$.
\end{lema}

\begin{proof} The ideal generator is a sum of operators, each acting on just one of the ten distinct marks,
so it suffices for each
factor to solve its own backward equation. For $O_{e_g}+\theta_g\chi_{e_g}$ we have
$\cL O_{e_g}=k_g\chi_{e_g}$ and
$\cL\chi_{e_g}=-4\alpha\chi_{e_g}$, and $\theta_g'+k_g-4\alpha\theta_g=0$ because
$\partial_t\Delta y=\partial_tz_t=-1/N$. For
$V_{d_i}-\xi_i\chi_{d_i}$ we have $\cL V_{d_i}=-\beta_i\chi_{d_i}$ and $\cL\chi_{d_i}=-\alpha\chi_{d_i}$,
and
$\xi_i'+\beta_i-\alpha\xi_i=0$. At $t=T$, $\Delta y=0$. The coefficient bounds follow from
$\pi_{e_g}=O(n^{-1})$ and $p_{d_i}=O(1)$.
\end{proof}

Multiplying each pair by $w(\bar r,\bar c)$ and summing gives the aggregate compensator at $t=0$, where
$\chi\equiv
V\equiv1$ and $O\equiv0$: \[
 \Phi_0=y_T^2\sum_{(e_1,e_2)}\pi_{e_1}\pi_{e_2}\bigl(1-y_Tp_{d_1}\bigr)\bigl(1-y_Tp_{d_2}\bigr)
 w(\bar r,\bar c),\qquad y_T=T/N .
\] Again, this is the solution of a linear system, not a factorization of the true law.

\subsection{Concentration}

\begin{lema}\label{lem:conc_inversa} For every fixed tolerance $\gamma>0$,
$\Pr\{|I_w(T)-\Phi_0|>\gamma n^2,\ T<\sigma'\}\le
e^{-c_{\zeta,\gamma,w}n}$, and a union bound over $O(n)$ horizons preserves the tail estimate. \end{lema}

\begin{proof} For $\ell=0,1,2$, the sum of the indicators with $\ell$ occupied edges is $O(n^{4-\ell})$.
There are at most $n$ choices
of queen per occupied edge and $n^2$ possible squares per available edge. In one step the sum of the
absolute changes is
$O(n^{3-\ell})$. A fixed diagonal meets any row or column in at most one square, affecting $O(n)$ pairs with
two occupied edges,
$O(n^2)$ with one, and $O(n^3)$ with none. For the indicators $\chi_{e_g}$ we use the fact that each line is
incident to $O(n)$
squares. Births give the same orders. The following table collects these counts together with the
coefficients of
Lemma~\ref{lem:compensador} and the remainders of the generator; for a birth we count the source state, with
$\ell-1$ occupied edges,
and not the resulting state.
\begin{center}\small
\begin{tabular}{@{}cccccc@{}}
\toprule
$\ell$ & patterns & coefficient & change per step & death remainder & birth remainder\\
\midrule
$0$ & $O(n^4)$ & $O(n^{-2})$ & $O(n^3)$ & $n^4\cdot n^{-2}\,(\eps/n+n^{-2})$ & ---\\
$1$ & $O(n^3)$ & $O(n^{-1})$ & $O(n^2)$ & $n^3\cdot n^{-1}\,(\eps/n+n^{-2})$
 & $n^4\cdot n^{-1}\,(\eps/n^2+n^{-3})$\\
$2$ & $O(n^2)$ & $O(1)$ & $O(n)$ & $n^2\cdot(\eps/n+n^{-2})$ & $n^3\cdot(\eps/n^2+n^{-3})$\\
\bottomrule
\end{tabular}
\end{center}
In each row, the coefficient times the change per step is $O(n)$, which bounds the increments, and each
remainder is $O(\eps n+1)$
per step. The second time differences and mixed terms sum to $O(1)$ per step by the derivative bounds in
Lemma~\ref{lem:compensador} and
the count of changes. We freeze both state and time at the stopping time $\sigma'$
(Lemma~\ref{lem:conc_directa}), including the exit
step. The probability of exit by $T_\zeta$ is $e^{-cn}$. The accumulated drift is $O_\zeta(\eps n^2+n)$, and
the martingale increments
are $O_\zeta(n)$ over $O(n)$ steps. Choosing the tolerance first and then $\eps$, Azuma's inequality gives
the bound. No union bound
over the $n^4$ pairs is needed. \end{proof}
